% paper/sections/03d_ridge_eikonal.tex
%
% §3.4  Ridge--Eikonal ハイブリッド再初期化
%   トポロジー変化を連続時間発展で扱う Gaussian 補助場 $\xiridge$ と
%   FMM/FSM による $\phi$ 再構成の二段構成．非一様格子拡張（D1--D4）まで含める．
%

% ═══════════════════════════════════════════════════════════════════════════════
\subsection{Ridge--Eikonal ハイブリッド再初期化}
\label{sec:ridge_eikonal}
% ═══════════════════════════════════════════════════════════════════════════════

\paragraph{本節の 2 用途．}
Ridge--Eikonal 経路は論文中で次の 2 つの目的で用いられる：
\textbf{(A) production stack の標準幾何更新．}~第~\ref{sec:algorithm}章の
7 段手順では，保存形移流後に $\psi=0.5$ 零等値面から FMM/FSM を起動して
$\phi$ を再構成し，$H_\varepsilon(-\phi)$ と質量閉鎖で CLS 変数へ戻す
（§~\ref{sec:ridge_fmm_reconstruction}）．縮約・診断経路では同じ操作を
$\psi$ 品質閾値に基づく幾何修復として起動してよい．
\textbf{(B) トポロジー変化の連続化．}~界面の融合・分裂時は
補助場 $\xiridge$ の時間発展（移流＋拡散）がリッジ集合の位相遷移を
連続変形として捉え，その後 FMM/FSM で $\phi$ を再構成する．
以下の前提・定式化は両用途に共通する．

\begin{tcolorbox}[defbox, title={本節の位置づけ}]
本節は，第~\ref{sec:levelset}章で導入した CLS 基本経路
（保存形移流・再初期化・$\psi$ 直接曲率）を\textbf{置き換える}ものではなく，
production stack の標準段階として，また品質劣化時とトポロジー変化時の
縮約・診断経路として $\phi$ を作り直す\textbf{幾何更新段}を与えるものである．
本章では連続定式化と再構成の論理に焦点を当て，
CCD/FCCD の離散演算子，非一様格子実装，感度較正の詳細は
第~\ref{sec:CCD}章，第~\ref{sec:grid}章，および
\S\ref{sec:ridge_eikonal_nu_grid} に委ねる．
\end{tcolorbox}

§~\ref{sec:reinit} で扱った CLS 圧縮--拡散平衡と，
古典 Level Set 文脈で用いられる Hamilton--Jacobi 型~\cite{JiangPeng2000}・FMM~\cite{Sethian1996}・FSM~\cite{Zhao2005} の諸手法は，
固定された零点集合を境界条件として Eikonal 条件 $|\bnabla\phi|=1$ を回復する枠組である．
そのままでは，界面の融合・分裂（coalescence / breakup）に必要な
「どの集合を新しい零点集合とみなすか」という自由度を与えない
（§~\ref{sec:notation} 囲みを参照）．
本節で用いる FMM/FSM の $\phi$ 再構成は
§~\ref{sec:ridge_fmm_reconstruction} で Ridge--Eikonal 用に改めて定式化する．
本節では\textbf{トポロジーの自由度}と\textbf{計量整合性}を
別々のスカラー場に担わせる二相スケジューリングを導入する．

\subsubsection{Gaussian 補助場 $\xiridge$ の定義（D0）}
\label{sec:ridge_aux_field}

各界面成分 $\Gamma_k$ との符号付き距離を $d_k(\bx,t)$ としたうえで，
\textbf{補助スカラー場} $\xiridge$ を以下で定義する：
\begin{equation}
  \xiridge(\bx,t) = \sum_k \exp\!\left(-\frac{d_k(\bx,t)^2}{\sigma_\xi^2}\right),
\end{equation}
ここで $\sigma_\xi$ は補助場の相互作用長であり，表面張力係数 $\sigma$ とは別の量である．
$\xiridge$ は\textbf{距離関数ではなく}
Eikonal 条件を満たさない；
これは局所極値・鞍点を許容するための意図的な設計であり，
正則なリッジ多様体と退化時刻を通じてトポロジー変化を表現する道具となる．

\paragraph{ブートストラップ初期化（$t=0$）．}
$d_k(\bx,0)$ は本スキームが時間発展中に計算する量と同じであるため，
$t=0$ では別経路で初期化する．具体的には以下の 2 つの方法のいずれかを採用する：
\begin{enumerate}[leftmargin=1.5em, itemsep=0pt, topsep=2pt]
  \item \textbf{幾何プリミティブ直接指定．}
        初期界面が解析形状（球・平面・楕円体等）で与えられる場合，
        各 $\Gamma_k$ の符号付き距離 $d_k(\bx,0)$ を解析式で直接評価する．
  \item \textbf{初期 $\psi$ からの FMM 再初期化．}
        初期 $\psi_0(\bx)$ が与えられる場合，
        §~\ref{sec:ridge_fmm_reconstruction} の FMM を
        零等値面 $\{\psi_0 = 0.5\}$ から起動し，$d_k(\bx,0)$ を構成する．
\end{enumerate}
いずれの場合も $\sigma_\xi$ は初期最小格子幅 $h_\text{min}(0)$ に対し
$\sigma_\xi \sim 3\,h_\text{min}(0)$ 程度に設定する．
一様初期格子では $h_\text{ref} = h_\text{min}$ が成り立つため，
D1 の $\sigma_\text{eff}$ はこの補助長の非一様格子版であり，
$\sigma_0/h_\text{ref} \equiv 3$ に対応する．

時間発展は移流・拡散で与える：
\begin{equation}
  \partial_t \xiridge + \bu\cdot\bnabla\xiridge
  = \varepsilon_\text{diff}\,\Delta \xiridge .
\end{equation}
拡散係数 $\varepsilon_\text{diff}$ は正則化（数値安定化）のためのものであり，
物理的拡散を表すわけではない．

\begin{remark}[他の補助場との区別]
$\xiridge$ は格子密度 Gaussian $\omegagrid(\phi) = \exp(-\alpha(\phi/\varepsilon_g)^2)$
（§~\ref{sec:grid_density}）とは\textbf{役割が直交}する：
$\omegagrid$ は格子間隔を制御し，$\xiridge$ は界面そのものを表す．
また，既存の \texttt{xi\_sdf} 再初期化モード
（\texttt{EikonalReinitializer.\_xi\_sdf\_phi}）が index 空間の
Euclidean 距離を計算して Eikonal を満たすのに対し，
$\xiridge$ は意図的に Eikonal を破る．
\end{remark}

\subsubsection{リッジ集合による界面定義}

界面 $\Gamma$ を $\xiridge$ の\textbf{リッジ集合}として定義する：
\begin{equation}
  \Gamma = \bigl\{\bx \,\big|\, \bnabla\xiridge(\bx) = 0,\;\;
  \lambda_{\min}\bigl(\bnabla^2\xiridge(\bx)\bigr) < 0 \bigr\},
  \label{eq:ridge_set_def}
\end{equation}
ここで $\lambda_{\min}(\cdot)$ は Hessian の最小固有値である．
単一の滑らかな界面では，$\xiridge$ は法線方向に極大，接線方向にほぼ平坦な
Morse--Bott 型のリッジ多様体を持つ．したがって式~\eqref{eq:ridge_set_def} は
孤立臨界点を列挙する通常の Morse 条件ではなく，法線 Hessian が負である
リッジ候補を抽出する条件として読む．
リッジ上の法線方向は $\lambda_{\min}$ に対応する単位固有ベクトル $\hat{\bm{v}}_{\min}$ で与え，
\begin{equation}
  \hat{\bm n}(\bx) =
  \begin{cases}
    \bnabla\xiridge(\bx)/|\bnabla\xiridge(\bx)|, & \bx \notin \Gamma, \\
    \hat{\bm{v}}_{\min}\bigl(\bnabla^2\xiridge(\bx)\bigr), & \bx \in \Gamma,
  \end{cases}
\end{equation}
と定める．$\Gamma$ の近傍で $\bnabla\xiridge\to 0$ のとき，法線方向の二次近似により
両定義は向きの符号選択を除いて同じ法線束へ収束する．
これは法線方向に極大・接線方向に連続というクレスト線（2次元）
／クレスト面（3次元）に対応する．
孤立界面ではリッジ集合は古典的界面と一致し，
複数界面の相互作用では\textbf{リッジ併合・分岐}が生じ，
これが融合・分裂に対応する．
したがってトポロジー変化は離散的なリメッシュ操作ではなく
\textbf{$\xiridge$ の連続変形}として表され，退化時刻だけを幾何イベントとして扱う．

\subsubsection{トポロジー変化の連続化}

従来の Level Set では，融合・分裂は $\phi$ の零等位集合のトポロジーが
不連続に変化する特異時刻として現れる．
Ridge--Eikonal 枠組では，これらの事象はすべて補助場 $\xiridge$ の
連続的な変形として単一パラメータ族内で表現される．

\begin{proposition}[リッジ集合の正則区間での連続性]
\label{prop:ridge_morse_continuity}
$\xiridge(\cdot, t)$ のリッジ集合が正則な Morse--Bott 多様体であり，
法線 Hessian の負固有値が 0 から離れている時間区間では，
$\Gamma(t)$ は以下の意味で連続変形する：
\begin{itemize}[leftmargin=1.5em, itemsep=0pt, topsep=2pt]
  \item 正則区間ではリッジの位置と法線束は時間に対して連続に変化する．
  \item 融合・分裂は，正則性が破れる孤立時刻をまたいでリッジ成分の接続関係が変わる事象として現れる．
  \item 退化候補は法線 Hessian の固有値と最小距離条件で検出し，FMM/FSM へ渡す零点集合を明示的に確定する．
\end{itemize}
\end{proposition}

\begin{proof}[\upshape 証明スケッチ]
$\xiridge \in C^2$ かつリッジ上の法線 Hessian が一様に非退化であれば，
陰関数定理により $\bnabla_\perp\xiridge=0$ の解多様体は時間に対して連続に追跡できる．
接線方向の零固有値は界面多様体そのものに対応するため，通常の Morse 臨界点として
孤立化する必要はない．
離散計算で非退化仮定が破れる場合の安全策は注記~\ref{rem:morse_degenerate_safeguard} を参照．
\end{proof}

\begin{remark}[退化リッジに対する数値的安全策]
\label{rem:morse_degenerate_safeguard}
命題~\ref{prop:ridge_morse_continuity} の法線非退化仮定は解析上は正則区間で自然だが，
離散 $\xiridge$ では法線 Hessian 固有値が 0 に近づく候補域が浮動小数演算で顕在化しうる．
命題~\ref{prop:ridge_morse_continuity} の仮定が離散計算で破れる可能性に対して
実装では再現性を壊さない\textbf{決定論的}な安全策を採るべきである：
\begin{enumerate}[leftmargin=1.5em, itemsep=0pt, topsep=2pt]
  \item \textbf{退化候補の flagging．}
        $|\lambda_{\perp,\min}(\bnabla^2\xiridge)| < \varepsilon_\text{ridge}$（推奨
        $\varepsilon_\text{ridge} = 10^{-6} \max|\bnabla^2\xiridge|$）の領域を，
        位相更新を保留すべき退化候補域として明示的に検出する．
  \item \textbf{局所平均による再評価．}
        隣接格子点間で $\lambda_{\min}$ の符号が反転する領域を
        「退化候補域」として flagging し，
        曲率評価では局所的に半径 $2h$ の球平均で Hessian を再構成して
        評価する．
\end{enumerate}
乱数摂動は結果の再現性を損ないうるため，本稿の主アルゴリズムでは仮定しない．
退化リッジの完全な分類は本稿の範囲外であり，ここでは
FMM/FSM へ渡す零点集合を決定論的に固定する実装契約だけを採用する．
\end{remark}

以上より明示的リメッシュは不要であり，
通常の界面移流で要求される CFL 制約の範囲で時間刻みを選べば，
Ridge 集合の連続変形は時間積分の中で自然に追跡される．
本稿では退化遷移専用の普遍的時定数を導入せず，
時間刻みの実務的選択は既存の移流・再初期化制約に従わせる．

\subsubsection{FMM/FSM による $\phi$ の再構成と一意性}
\label{sec:ridge_fmm_reconstruction}

\begin{theorem}[FMM/FSM による符号付き $\phi$ 一意再構成]
\label{thm:eikonal_reconstruction}
リッジ集合 $\Gamma$ と相ラベル（気相側・液相側を指定する符号）が以下を満たすとき：
\begin{itemize}[leftmargin=1.5em, itemsep=0pt, topsep=2pt]
  \item\textbf{幾何正則性．}$\Gamma$ は閉じた自己交差しない $(d-1)$ 次元 Lipschitz 多様体．
  \item\textbf{射影一意性．}$\Gamma$ の近傍点からの最近接射影が一意である．
  \item\textbf{符号一貫性．}気相側を正，液相側を負とする相ラベルが $\Gamma$ の両側で連続に定まる．
\end{itemize}
（数値検証の詳細は §~\ref{sec:ridge_admissibility}）．Eikonal 再構成問題
\begin{equation}
  |\bnabla\phi| = 1 \text{ in } \Omega \setminus \Gamma, \qquad
  \phi = 0 \text{ on } \Gamma,
\end{equation}
には，CLS 規約（気相側で正・液相側で負；§~\ref{sec:notation}）に従う
唯一の符号付き粘性解 $\phi(\bx) = s(\bx)\,d(\bx, \Gamma)$ が存在する．
ここで $s(\bx)\in\{+1,-1\}$ は相ラベルで定まる符号である．
\end{theorem}

\begin{proof}[\upshape 証明スケッチ（Crandall--Lions 1983~\cite{CrandallLions1983}）]
以下，符号付き距離（気相側で正）を $\phi=s\,d(\bx,\Gamma)$ と書く．
(i)~\textbf{存在．} $\Gamma$ が閉集合なら $d(\bx,\Gamma)$ は Lipschitz 連続である．
(ii)~\textbf{局所微分可能性．} 最近接射影 $\pi : \bx\mapsto \Gamma$ が一意な領域で
$\phi$ は微分可能で $\bnabla\phi(\bx) = \hat{\bm n}_\Gamma(\pi(\bx))$．
(iii)~\textbf{Eikonal 成立．} したがって $|\bnabla\phi|=1$ がほぼ至るところで成立．
(iv)~\textbf{一意性．} 粘性解の比較原理~\cite{CrandallLions1983} により
$\phi$ は $\Gamma$ を境界条件として一意．
\end{proof}

\paragraph{アルゴリズム}
$\phi$ の再構成は FMM（狭帯域・非一様推奨）または FSM（全域）を適用する．
両者で得られる $\phi$ は粘性解として同一である．
ここで Dirichlet 条件として与える $\Gamma$ は，
固定トポロジーの再初期化では入力 $\phi_0$ の零点集合
（または $\psi=0.5$ の線形補間零交差集合）そのものでなければならない．
Gaussian 補助場 $\xiridge$ のリッジ集合は，
トポロジー変化を明示的に扱う補助経路で界面候補を定義するための場であり，
通常の固定トポロジー再初期化において
``$\xiridge$ の局所極大かつ $|\phi_0|<O(h)$'' という判定だけで
新たな $\phi=0$ Dirichlet 点を追加してはならない．
そのような追加は定理~\ref{thm:eikonal_reconstruction} の境界集合を
$\Gamma$ から $\Gamma\cup\Gamma_{\mathrm{art}}$ へ変更し，
零点集合保存性と周期・反射対称性を同時に破る別問題を解くことになる．
再構成後は $\psi = H_{\varepsilon_\text{local}}(-\phi)$ に戻して
標準の CLS パイプラインへ復帰する：
すなわち，第~\ref{sec:algo_flow}節の第1段では
通常どおり TVD-RK3 + FCCD 保存形面フラックスによる $\psi$ 移流を継続し，
必要に応じて後段で FCCD 系の面中心演算子を
圧力・CSF・運動量離散化に接続する．

\subsubsection{幾何的許容性条件}
\label{sec:ridge_admissibility}

定理~\ref{thm:eikonal_reconstruction} の条件（幾何正則性・射影一意性・符号一貫性）に加え，
数値実装では以下の 1 条件が必要となる：
\begin{enumerate}
  \item\textbf{法線一貫性．}$\xiridge$ から推定される法線方向が $\Gamma$ 上で整合（向き反転なし）．
\end{enumerate}
幾何正則性と射影一意性は，曲率上限 $\kappa_{\max}\lesssim 1/h$ と最小間隔
$d_{\min}\gtrsim 2\varepsilon$ を満たすかをリッジ抽出時に確認する．
法線一貫性は Hessian 固有ベクトル符号を連続 branch で選択することで確保する．

\subsubsection{表面張力ありの界面幅整合}

純粋 FMM 再構成は，表面張力が有意な場合には
SDF 幅と diffuse-interface 幅の不一致に敏感になりやすい．
根本原因は，PPE 側で参照される表面張力残差が
$\sigma\kappa/\varepsilon$（ここで $\sigma$ は表面張力係数）にスケールする一方，
再構成された SDF が局所界面幅を必ずしも保持しない点にある．
\textbf{$\varepsilon$-widening}（式~\eqref{eq:eps_local}）を用いる場合は
この不一致を補正するための明示的な感度・後処理条件として登録し，
標準経路の隠れた正則化として扱わない．

\subsubsection{面中心追跡との接続}
\label{sec:ridge_fccd_handoff}

トポロジーが固定された段階では，
再構成された $\phi$ から $\psi = H_{\varepsilon_\text{local}}(-\phi)$ を再生成し，
以降は通常の CLS 更新へ戻る．
したがって本稿の既定経路では，
界面追跡そのものは第~\ref{sec:advection}章で定義する
TVD-RK3 + FCCD 保存形面フラックスの $\psi$ 移流が担い，
FCCD 面中心演算子は後段で圧力・CSF・運動量離散化に接続される．
退化リッジ候補を含まない時間帯（トポロジー変化が発生しない区間）では，
アルゴリズムは
(i) 標準 CLS $\psi$ 移流，(ii) production では各 step の Ridge--Eikonal 再投影，
縮約・診断経路では必要時の Ridge--Eikonal 再投影
という同じ幾何更新契約として読める．

\paragraph{本節の要点}
\begin{itemize}
  \item $\xiridge$ と $\phi$ は\textbf{役割が相補}：
        前者はトポロジー自由度，後者は計量整合性を担う．
  \item トポロジー変化は $\xiridge$ の連続変形と退化候補検出として扱う．
  \item 一意性は粘性解の比較原理により保証される．
  \item 表面張力ありで $\varepsilon$-widening を用いる場合は，標準経路と分けて明示する．
\end{itemize}
